\chapter{2022年新高考I卷}

\section{单选题}

\begin{problem}
若集合 \(M=\{x\mid \sqrt{x}<4\}\)，\(N=\{x\mid 3x\ge 1\}\)，则 \(M\cap N=\)（\quad）

\choices
    {\(\{x\mid 0\le x<2\}\)}
    {\(\left\{x\mid \frac13\le x<2\right\}\)}
    {\(\{x\mid 3\le x<16\}\)}
    {\(\left\{x\mid \frac13\le x<16\right\}\)}
\end{problem}

\begin{answer}
D
\end{answer}

\begin{solution}
因为 \(\sqrt{x}\) 有意义要求 \(x\ge 0\)，而 \(\sqrt{x}<4\) 两边均为非负数，平方得 \(x<16\)，所以 \(M=\left\{x\mid 0\le x<16\right\}\). 由 \(3x\ge 1\) 得 \(x\ge \frac13\)，所以 \(N=\left\{x\mid x\ge \frac13\right\}\). 因此 \(M\cap N=\left\{x\mid \frac13\le x<16\right\}\). 故选 D.
\end{solution}

\begin{problem}
若 \(\i(1-z)=1\)，则 \(z+\overline{z}=\)（\quad）

\choices
    {\(-2\)}
    {\(-1\)}
    {1}
    {2}
\end{problem}

\begin{answer}
D
\end{answer}

\begin{solution}
因为 \(\i\left(1-z\right)=1\)，所以 \(1-z=\frac{1}{\i}=-\i\)，这里用到了 \(\i\cdot \left(-\i\right)=1\). 于是 \(z=1+\i\)，\(\overline{z}=1-\i\). 相加得 \(z+\overline{z}=2\). 故选 D.
\end{solution}

\begin{problem}
在 \(\triangle ABC\) 中，点 \(D\) 在边 \(AB\) 上，\(BD=2DA\). 记 \(\overrightarrow{CA}=\bs m\)，\(\overrightarrow{CD}=\bs n\)，则 \(\overrightarrow{CB}=\)（\quad）

\choices
    {\(3\bs m-2\bs n\)}
    {\(-2\bs m+3\bs n\)}
    {\(3\bs m+2\bs n\)}
    {\(2\bs m+3\bs n\)}
\end{problem}

\begin{answer}
B
\end{answer}

\begin{solution}
\FigureLayoutDeclare{img/2022/new_gaokao_paper_1/q3_solution_fig1.png}{5.0cm}{0bp 0bp 0bp 0bp}
\solutionfigure{\bitmapfigure[max width=0.22\textwidth]{img/2022/new_gaokao_paper_1/q3_solution_fig1.png}}

因为点 \(D\) 在边 \(AB\) 上且 \(BD=2DA\)，所以 \(\overrightarrow{AD}=\frac13\overrightarrow{AB}\). 于是 \(\bs n=\overrightarrow{CD}=\overrightarrow{CA}+\overrightarrow{AD}=\bs m+\frac13\overrightarrow{AB}\). 又 \(\overrightarrow{AB}=\overrightarrow{CB}-\overrightarrow{CA}=\overrightarrow{CB}-\bs m\)，代入得 \(\bs n=\bs m+\frac13\left(\overrightarrow{CB}-\bs m\right)=\frac23\bs m+\frac13\overrightarrow{CB}\). 两边乘以 \(3\) 并移项得 \(\overrightarrow{CB}=3\bs n-2\bs m=-2\bs m+3\bs n\). 故选 B.
\end{solution}

\begin{problem}
南水北调工程缓解了北方一些地区水资源短缺问题，其中一部分水蓄入某水库. 已知该水库水位为海拔 \(148.5\text{m}\) 时，相应水面的面积为 \(140.0\text{km}^2\)；水位为海拔 \(157.5\text{m}\) 时，相应水面的面积为 \(180.0\text{km}^2\). 将该水库在这两个水位间的形状看作一个棱台，则该水库水位从海拔 \(148.5\text{m}\) 上升到 \(157.5\text{m}\) 时，增加的水量约为（\(\sqrt7\approx 2.65\)）（\quad）

\choices
    {\(1.0\times 10^9\text{m}^3\)}
    {\(1.2\times 10^9\text{m}^3\)}
    {\(1.4\times 10^9\text{m}^3\)}
    {\(1.6\times 10^9\text{m}^3\)}
\end{problem}

\begin{answer}
C
\end{answer}

\begin{solution}
因为水位从海拔 \(148.5\text{ m}\) 上升到 \(157.5\text{ m}\)，所以棱台的高为 \(h=157.5-148.5=9\text{ m}\). 记下底面积为 \(S=140.0\text{ km}^2=140\times 10^6\text{ m}^2\)，上底面积为 \(S^{\prime}=180.0\text{ km}^2=180\times 10^6\text{ m}^2\)，这里用到了 \(1\text{ km}^2=10^6\text{ m}^2\). 由棱台体积公式 \(V=\frac13h\left(S+S^{\prime}+\sqrt{SS^{\prime}}\right)\)，而 \(\sqrt{SS^{\prime}}=\sqrt{140\times 180}\times 10^6\text{ m}^2=60\sqrt{7}\times 10^6\text{ m}^2\approx 60\times 2.65\times 10^6\text{ m}^2=159\times 10^6\text{ m}^2\). 所以 \(V\approx \frac13\times 9\times \left(140+180+159\right)\times 10^6\text{ m}^3=1437\times 10^6\text{ m}^3=1.437\times 10^9\text{ m}^3\). 故增加的水量约为 \(1.4\times 10^9\text{ m}^3\)，选 C.
\end{solution}

\begin{problem}
从 \(2\) 至 \(8\) 的 \(7\) 个整数中随机取 \(2\) 个不同的数，则这 \(2\) 个数互质的概率为（\quad）

\choices
    {\(\dfrac16\)}
    {\(\dfrac13\)}
    {\(\dfrac12\)}
    {\(\dfrac23\)}
\end{problem}

\begin{answer}
D
\end{answer}

\begin{solution}
从 \(2\) 至 \(8\) 的 \(7\) 个整数中任取 \(2\) 个不同数的取法共有 \(\mathrm C_7^2=21\) 种. 若两数不互质，则它们有大于 \(1\) 的公约数. 在 \(2\)，\(3\)，\(4\)，\(5\)，\(6\)，\(7\)，\(8\) 中，偶数为 \(2\)，\(4\)，\(6\)，\(8\)，其中任取 \(2\) 个均不互质，共有 \(\mathrm C_4^2=6\) 种，即 \(\left(2,4\right)\)，\(\left(2,6\right)\)，\(\left(2,8\right)\)，\(\left(4,6\right)\)，\(\left(4,8\right)\)，\(\left(6,8\right)\). 既非全偶又不互质的只能是 \(3\) 与 \(6\) 的组合，即 \(\left(3,6\right)\). 其余含 \(5\)，\(7\) 或 \(3\) 与非 \(6\) 偶数的组合均互质. 所以不互质的取法共 \(6+1=7\) 种，互质的取法共 \(21-7=14\) 种. 因此所求概率为 \(\frac{14}{21}=\frac23\). 故选 D.
\end{solution}

\begin{problem}
记函数 \(f(x)=\sin\left(\omega x+\frac{\pi}{4}\right)+b\)（\(\omega>0\)） 的最小正周期为 \(T\). 若 \(\frac{2\pi}{3}<T<\pi\)，且 \(y=f(x)\) 的图象关于点 \(\left(\frac{3\pi}{2},2\right)\) 中心对称，则 \(f\left(\frac{\pi}{2}\right)=\)（\quad）

\choices
    {1}
    {\(\dfrac32\)}
    {\(\dfrac52\)}
    {3}
\end{problem}

\begin{answer}
A
\end{answer}

\begin{solution}
因为 \(f\left(x\right)=\sin\left(\omega x+\frac{\pi}{4}\right)+b\) 的最小正周期为 \(T=\frac{2\pi}{\omega}\)，由 \(\frac{2\pi}{3}<T<\pi\) 得 \(2<\omega<3\). 若 \(y=f\left(x\right)\) 的图象关于点 \(\left(\frac{3\pi}{2},2\right)\) 中心对称，则正弦部分的中心纵坐标为 \(0\)，所以 \(b=2\)，且 \(\sin\left(\frac{3\pi}{2}\omega+\frac{\pi}{4}\right)=0\). 于是存在 \(k\in\Z\)，使得 \(\frac{3\pi}{2}\omega+\frac{\pi}{4}=k\pi\). 两边除以 \(\pi\) 得 \(\frac32\omega+\frac14=k\)，即 \(\omega=\frac{4k-1}{6}\). 结合 \(2<\omega<3\) 得 \(2<\frac{4k-1}{6}<3\)，即 \(\frac{13}{4}<k<\frac{19}{4}\). 因为 \(k\in\Z\)，所以 \(k=4\)，从而 \(\omega=\frac{15}{6}=\frac52\). 因此 \(f\left(x\right)=\sin\left(\frac52x+\frac{\pi}{4}\right)+2\)，所以 \(f\left(\frac{\pi}{2}\right)=\sin\left(\frac{5\pi}{4}+\frac{\pi}{4}\right)+2=\sin\left(\frac{3\pi}{2}\right)+2=-1+2=1\). 故选 A.
\end{solution}

\begin{problem}
设 \(a=0.1\e^{0.1}\)，\(b=\frac{0.1}{1-0.1}\)，\(c=-\ln 0.9\)，则（\quad）

\choices
    {\(a<b<c\)}
    {\(c<b<a\)}
    {\(c<a<b\)}
    {\(a<c<b\)}
\end{problem}

\begin{answer}
C
\end{answer}

\begin{solution}
记 \(x=0.1\)，则 \(a=x\e^{x}\)，\(b=\frac{x}{1-x}=\frac19\)，\(c=-\ln\left(1-x\right)\). 先证 \(a<b\). 设 \(h\left(x\right)=\left(1-x\right)\e^{x}\)，\(x\in\left[0,0.1\right]\)，则 \(h^{\prime}\left(x\right)=-\e^{x}+\left(1-x\right)\e^{x}=-x\e^{x}\le 0\)，且仅在 \(x=0\) 处取零，所以 \(h\left(x\right)\) 在 \(\left[0,0.1\right]\) 上单调递减. 因为 \(h\left(0\right)=1\)，所以 \(h\left(0.1\right)<1\)，即 \(0.9\e^{0.1}<1\)，也就是 \(\e^{0.1}<\frac{10}{9}\). 两边乘以 \(0.1\) 得 \(0.1\e^{0.1}<\frac19\)，即 \(a<b\). 再证 \(c<a\). 设 \(k\left(x\right)=x\e^{x}+\ln\left(1-x\right)\)，\(x\in\left[0,0.1\right]\)，则 \(k\left(0\right)=0\)，且 \(k^{\prime}\left(x\right)=\e^{x}+x\e^{x}-\frac{1}{1-x}=\left(1+x\right)\e^{x}-\frac{1}{1-x}\)，\(k^{\prime}\left(0\right)=0\). 又 \(k^{\prime\prime}\left(x\right)=\left(2+x\right)\e^{x}-\frac{1}{\left(1-x\right)^2}\). 为证 \(k^{\prime\prime}\left(x\right)>0\)，先证 \(\e^{x}\ge 1+x\). 设 \(s\left(x\right)=\e^{x}-x-1\)，则 \(s^{\prime}\left(x\right)=\e^{x}-1\ge 0\)，\(x\ge 0\)，所以 \(s\left(x\right)\) 单调递增，又 \(s\left(0\right)=0\)，故 \(x\ge 0\) 时 \(\e^{x}\ge 1+x\). 于是 \(\left(2+x\right)\e^{x}\ge \left(2+x\right)\left(1+x\right)\). 而 \(\left(2+x\right)\left(1+x\right)-\frac{1}{\left(1-x\right)^2}=\frac{1-x-3x^2+x^3+x^4}{\left(1-x\right)^2}\)，分子 \(=1-x-3x^2+x^3+x^4\ge 1-0.1-0.03=0.87>0\)，\(x\in\left[0,0.1\right]\)，所以 \(k^{\prime\prime}\left(x\right)>0\)，\(x\in\left[0,0.1\right]\). 因此 \(k^{\prime}\left(x\right)\) 单调递增，又 \(k^{\prime}\left(0\right)=0\)，故 \(x\in\left(0,0.1\right]\) 时 \(k^{\prime}\left(x\right)>0\). 从而 \(k\left(x\right)\) 单调递增，又 \(k\left(0\right)=0\)，故 \(k\left(0.1\right)>0\)，即 \(0.1\e^{0.1}+\ln 0.9>0\)，也就是 \(-\ln 0.9<0.1\e^{0.1}\)，即 \(c<a\). 结合 \(a<b\) 得 \(c<a<b\). 故选 C.
\end{solution}

\begin{problem}
已知正四棱锥的侧棱长为 \(l\)，其各顶点都在同一球面上. 若该球的体积为 \(36\pi\)，且 \(3\le l\le 3\sqrt3\)，则该正四棱锥体积的取值范围是（\quad）

\choices
    {\(\left[18,\frac{81}{4}\right]\)}
    {\(\left[\frac{27}{4},\frac{81}{4}\right]\)}
    {\(\left[\frac{27}{4},\frac{64}{3}\right]\)}
    {\([18,27]\)}
\end{problem}

\begin{answer}
C
\end{answer}

\begin{solution}
\solutionfigure{\bitmapfigure[max width=0.30\textwidth]{img/2022/new_gaokao_paper_1/q8_solution_fig1.png}}

因为球体积 \(\frac43\pi R^3=36\pi\)，所以 \(R^3=27\)，故 \(R=3\). 设正四棱锥底面正方形边长为 \(2a\)，\(a>0\)，高为 \(h\)，\(h>0\)，侧棱长为 \(l\)，则底面对角线长的一半为 \(\sqrt{2}a\)，故 \(l^2=2a^2+h^2\). 又各顶点共球，球心在高线上，所以 \(R^2=2a^2+\left(h-R\right)^2\)，即 \(2a^2=2Rh-h^2\). 代入 \(R=3\) 得 \(2a^2=6h-h^2\). 与 \(l^2=2a^2+h^2\) 联立消去 \(a\) 得 \(l^2=6h\)，即 \(h=\frac{l^2}{6}\). 体积为 \(V=\frac13\cdot \left(2a\right)^2h=\frac43a^2h=\frac23h\left(l^2-h^2\right)\). 代入 \(h=\frac{l^2}{6}\) 得 \(V\left(l\right)=\frac{l^4}{9}-\frac{l^6}{324}\)，\(3\le l\le 3\sqrt{3}\). 求导得 \(V^{\prime}\left(l\right)=\frac{4l^3}{9}-\frac{l^5}{54}=\frac{l^3\left(24-l^2\right)}{54}\). 令 \(V^{\prime}\left(l\right)=0\) 得 \(l=2\sqrt{6}\)，且 \(3<2\sqrt{6}<3\sqrt{3}\). 当 \(3\le l<2\sqrt{6}\) 时 \(V^{\prime}\left(l\right)>0\)，当 \(2\sqrt{6}<l\le 3\sqrt{3}\) 时 \(V^{\prime}\left(l\right)<0\)，所以 \(V\left(l\right)\) 先增后减，最大值在 \(l=2\sqrt{6}\) 处取得，为 \(V\left(2\sqrt{6}\right)=\frac{576}{9}-\frac{13824}{324}=64-\frac{128}{3}=\frac{64}{3}\). 端点值为 \(V\left(3\right)=9-\frac{729}{324}=\frac{27}{4}\)，\(V\left(3\sqrt{3}\right)=81-\frac{19683}{324}=\frac{81}{4}\). 因为 \(\frac{27}{4}<\frac{81}{4}\)，所以最小值为 \(\frac{27}{4}\). 故体积取值范围是 \(\left[\frac{27}{4},\frac{64}{3}\right]\)，选 C.
\end{solution}
\section{多选题}

\begin{problem}
已知正方体 \(ABCD-A_1B_1C_1D_1\)，则（\quad）

\choices
    {直线 \(BC_1\) 与 \(DA_1\) 所成的角为 \(90^\circ\)}
    {直线 \(BC_1\) 与 \(CA_1\) 所成的角为 \(90^\circ\)}
    {直线 \(BC_1\) 与平面 \(BB_1D_1D\) 所成的角为 \(45^\circ\)}
    {直线 \(BC_1\) 与平面 \(ABCD\) 所成的角为 \(45^\circ\)}
\end{problem}

\begin{answer}
ABD
\end{answer}

\begin{solution}
\solutionfigure{\bitmapfigure[max width=0.25\textwidth]{img/2022/new_gaokao_paper_1/q9_solution_fig1.png}}

设正方体棱长为 \(1\)，取 \(A(0,0,0)\)，\(B(1,0,0)\)，\(C(1,1,0)\)，\(D(0,1,0)\)，\(A_1(0,0,1)\)，\(B_1(1,0,1)\)，\(C_1(1,1,1)\). 因为 \(\overrightarrow{BC_1}=(0,1,1)\)，\(\overrightarrow{DA_1}=(0,-1,1)\)，内积为 \(0\)，所以两直线垂直，故A正确. 因为 \(\overrightarrow{CA_1}=(-1,-1,1)\) 与 \(\overrightarrow{BC_1}\) 的内积为 \(0\)，所以 \(BC_1\perp CA_1\)，故B正确. 因为 \(BB_1\perp A_1C_1\) 且 \(B_1D_1\perp A_1C_1\)，所以 \(A_1C_1\perp\) 平面 \(BB_1D_1D\). 设 \(O_1\) 为 \(A_1C_1\) 与 \(B_1D_1\) 的交点，则 \(C_1O_1\perp\) 平面 \(BB_1D_1D\)，而 \(B\) 在该平面内，所以 \(BO_1\) 为 \(BC_1\) 在该平面上的射影，所成角为 \(\angle C_1BO_1\). 因为 \(\lvert C_1O_1\rvert=\frac{\sqrt{2}}{2}\)，\(\lvert BC_1\rvert=\sqrt{2}\)，所以 \(\sin\angle C_1BO_1=\frac{1}{2}\)，夹角为 \(30^\circ\)，故C错误. 因为 \(C_1C\perp\) 平面 \(ABCD\)，所以 \(BC_1\) 与底面所成角为 \(\angle C_1BC=45^\circ\)，故D正确. 综上选ABD.
\end{solution}

\begin{problem}
已知函数 \(f(x)=x^3-x+1\)，则（\quad）

\choices
    {\(f(x)\) 有两个极值点}
    {\(f(x)\) 有三个零点}
    {点 \((0,1)\) 是曲线 \(y=f(x)\) 的对称中心}
    {直线 \(y=2x\) 是曲线 \(y=f(x)\) 的切线}
\end{problem}

\begin{answer}
AC
\end{answer}

\begin{solution}
因为 \(f'(x)=3x^2-1\)，令 \(f'(x)=0\) 得 \(x=\pm\frac{\sqrt{3}}{3}\)，且 \(f'\) 在这两点两侧变号，所以 \(f(x)\) 有两个极值点，故A正确. 因为当 \(x<-\frac{\sqrt{3}}{3}\) 或 \(x>\frac{\sqrt{3}}{3}\) 时 \(f'(x)>0\)，当 \(-\frac{\sqrt{3}}{3}<x<\frac{\sqrt{3}}{3}\) 时 \(f'(x)<0\)，所以 \(x=-\frac{\sqrt{3}}{3}\) 为极大值点，\(x=\frac{\sqrt{3}}{3}\) 为极小值点. 又 \(f\left(-\frac{\sqrt{3}}{3}\right)=1+\frac{2\sqrt{3}}{9}>0\)，\(f\left(\frac{\sqrt{3}}{3}\right)=1-\frac{2\sqrt{3}}{9}>0\)，且当 \(x\to -\infty\) 时 \(f(x)\to -\infty\)，当 \(x\to +\infty\) 时 \(f(x)\to +\infty\)，所以 \(f(x)\) 仅在 \(\left(-\infty,-\frac{\sqrt{3}}{3}\right)\) 内有一个零点，故B错误. 因为 \(h(x)=x^3-x\) 满足 \(h(-x)=-h(x)\)，为奇函数，图象关于原点对称，而 \(f(x)=h(x)+1\) 为其向上平移 \(1\) 个单位，所以 \(y=f(x)\) 关于 \((0,1)\) 中心对称，故C正确. 因为令 \(f'(x)=2\) 得 \(x=\pm 1\)，\(f(1)=1\) 处切线为 \(y=2x-1\)，\(f(-1)=1\) 处切线为 \(y=2x+3\)，斜率为 \(2\) 的切点仅此两处，均不给出 \(y=2x\)，所以 \(y=2x\) 不是切线，故D错误. 综上选AC.
\end{solution}

\begin{problem}
已知 \(O\) 为坐标原点，点 \(A(1,1)\) 在抛物线 \(C:\ x^2=2py\)（\(p>0\)） 上，过点 \(B(0,-1)\) 的直线交 \(C\) 于 \(P\)，\(Q\) 两点，则（\quad）

\choices
    {\(C\) 的准线为 \(y=-1\)}
    {直线 \(AB\) 与 \(C\) 相切}
    {\(|OP|\cdot |OQ|>|OA|^2\)}
    {\(|BP|\cdot |BQ|>|BA|^2\)}
\end{problem}

\begin{answer}
BCD
\end{answer}

\begin{solution}
由 \(A(1,1)\) 在 \(x^2=2py\) 上得 \(1=2p\)，即 \(p=\frac{1}{2}\)，故抛物线为 \(x^2=y\)，准线为 \(y=-\frac{1}{4}\)，所以A错误. 因为 \(k_{AB}=\frac{1-(-1)}{1-0}=2\)，直线 \(AB\) 为 \(y=2x-1\)，联立 \(y=x^2\) 得 \(x^2-2x+1=0\)，即 \((x-1)^2=0\)，有重根 \(x=1\)，所以 \(AB\) 与 \(C\) 相切，故B正确. 设过 \(B(0,-1)\) 交 \(C\) 于 \(P\)，\(Q\) 两点的直线为 \(l\)，记 \(P(x_1,y_1)\)，\(Q(x_2,y_2)\). 若 \(l\) 斜率不存在，则 \(l\) 为 \(x=0\) 或 \(x=t\ne 0\)，每条竖直线与 \(y=x^2\) 至多交于一点，不构成两点 \(P\)，\(Q\)，所以交于两点的 \(l\) 必有斜率，设为 \(y=kx-1\). 联立 \(x^2=kx-1\) 得 \(x^2-kx+1=0\)，有两个不同交点的充要条件是 \(\Delta=k^2-4>0\)，即 \(\lvert k\rvert>2\)，此时根与系数关系为 \(x_1+x_2=k\)，\(x_1x_2=1\)，且 \(y_1=x_1^2\)，\(y_2=x_2^2\). 因为 \(\lvert OP\rvert=\sqrt{x_1^2+y_1^2}=\lvert x_1\rvert\sqrt{1+x_1^2}\)，\(\lvert OQ\rvert=\lvert x_2\rvert\sqrt{1+x_2^2}\)，所以 \(\lvert OP\rvert\cdot\lvert OQ\rvert=\lvert x_1x_2\rvert\sqrt{(1+x_1^2)(1+x_2^2)}=\sqrt{1+x_1^2+x_2^2+x_1^2x_2^2}\). 又 \(x_1^2+x_2^2=(x_1+x_2)^2-2x_1x_2=k^2-2\)，\(x_1^2x_2^2=1\)，所以上式等于 \(\sqrt{2+k^2-2}=\lvert k\rvert\). 因为 \(\lvert k\rvert>2\)，所以 \(\lvert OP\rvert\cdot\lvert OQ\rvert=\lvert k\rvert>2\). 又 \(\lvert OA\rvert^2=1^2+1^2=2\)，所以 \(\lvert OP\rvert\cdot\lvert OQ\rvert>\lvert OA\rvert^2\)，故C正确. 因为 \(y_1+1=kx_1\)，\(y_2+1=kx_2\)，所以 \(\lvert BP\rvert=\sqrt{x_1^2+(y_1+1)^2}=\lvert x_1\rvert\sqrt{1+k^2}\)由相同计算 \(\lvert BQ\rvert=\lvert x_2\rvert\sqrt{1+k^2}\)，于是 \(\lvert BP\rvert\cdot\lvert BQ\rvert=\lvert x_1x_2\rvert(1+k^2)=1+k^2\). 又 \(\lvert BA\rvert^2=(1-0)^2+(1+1)^2=5\)，而 \(1+k^2>5\) 等价于 \(k^2>4\)，恰为 \(\Delta>0\) 的条件，所以 \(\lvert BP\rvert\cdot\lvert BQ\rvert>\lvert BA\rvert^2\)，故D正确. 综上选BCD.
\end{solution}

\begin{problem}
已知函数 \(f(x)\) 及其导函数 \(f'(x)\) 的定义域均为 \(\R\)，记 \(g(x)=f'(x)\). 若 \(f\left(\frac32-2x\right)\)，\(g(2+x)\) 均为偶函数，则（\quad）

\choices
    {\(f(0)=0\)}
    {\(g\left(-\frac12\right)=0\)}
    {\(f(-1)=f(4)\)}
    {\(g(-1)=g(2)\)}
\end{problem}

\begin{answer}
BC
\end{answer}

\begin{solution}
因为 \(f\left(\frac{3}{2}-2x\right)\) 为偶函数，所以对一切 \(x\in\R\) 有 \(f\left(\frac{3}{2}+2x\right)=f\left(\frac{3}{2}-2x\right)\). 令 \(t=2x\)，则 \(f\left(\frac{3}{2}+t\right)=f\left(\frac{3}{2}-t\right)\)，即 \(f(x)\) 关于 \(x=\frac{3}{2}\) 对称，故 \(f(3-x)=f(x)\). 于是 \(f(-1)=f(3-(-1))=f(4)\)，所以C正确. 因为 \(g(2+x)\) 为偶函数，所以 \(g(2-x)=g(2+x)\)，即 \(g(4-x)=g(x)\)，故 \(g(x)\) 关于 \(x=2\) 对称. 对 \(f(3-x)=f(x)\) 两边求导得 \(-f'(3-x)=f'(x)\)，即 \(g(3-x)=-g(x)\). 由 \(g(4-x)=g(x)\) 得 \(g(x+1)=g(4-(x+1))=g(3-x)=-g(x)\)，所以 \(g(x+1)=-g(x)\)，进而 \(g(x+2)=-g(x+1)=g(x)\)，即 \(g\) 以 \(2\) 为周期. 在 \(g(3-x)=-g(x)\) 中令 \(x=\frac{3}{2}\) 得 \(g\left(\frac{3}{2}\right)=0\)，再由周期 \(2\) 得 \(g\left(-\frac{1}{2}\right)=g\left(\frac{3}{2}\right)=0\)，所以B正确. 因为 \(f(0)\) 的值可平移而不破坏对称关系，不能由题设唯一确定，例如取 \(g(x)=\cos(\pi x)\) 及其原函数 \(f(x)=\frac{\sin(\pi x)}{\pi}+C\) 即满足全部对称条件，而 \(C\) 任意，所以A不能成立. 因为由周期与 \(g(x+1)=-g(x)\) 得 \(g(-1)=g(1)\) 且 \(g(2)=-g(1)\)，所以 \(g(-1)=g(2)\) 当且仅当 \(g(2)=0\)，而上例中 \(g(2)=1\ne 0\)，此时 \(g(-1)=-1\ne g(2)\)，故D不成立. 综上选BC.
\end{solution}
\section{填空题}

\begin{problem}
\( \left(1-\frac{y}{x}\right)(x+y)^8 \) 的展开式中 \(x^2y^6\) 的系数为\fillinblank{}（用数字作答）.
\end{problem}

\begin{answer}
\(-28\)
\end{answer}

\begin{solution}
因为 \(\left(1-\frac{y}{x}\right)(x+y)^8=(x+y)^8-\frac{y}{x}(x+y)^8\). 在 \((x+y)^8\) 中，\(x^2y^6\) 来自 \(k=6\) 项，系数为 \(\mathrm C_8^6=28\). 在 \(\frac{y}{x}(x+y)^8\) 中，要得到 \(x^2y^6\)，需 \((x+y)^8\) 提供 \(x^3y^5\)，即 \(k=5\) 项，系数为 \(\mathrm C_8^5=56\)，带负号后贡献为 \(-56\). 所以所求系数为 \(28-56=-28\). 故答案为 \(-28\).
\end{solution}

\begin{problem}
写出与圆 \(x^2+y^2=1\) 和 \((x-3)^2+(y-4)^2=16\) 都相切的一条直线的方程\fillinblank{}.
\end{problem}

\begin{answer}
\(x=-1\)
\end{answer}

\begin{solution}
\solutionfigure{\bitmapfigure[max width=0.32\textwidth]{img/2022/new_gaokao_paper_1/q14_solution_fig1.png}}

记两圆圆心为 \(O(0,0)\)，\(O_1(3,4)\)，半径分别为 \(1\)，\(4\). 因为 \(\lvert OO_1\rvert=5=1+4\)，所以两圆外切，共有三条公切线. 直线 \(x=-1\) 到 \(O\) 的距离为 \(1\)，到 \(O_1\) 的距离为 \(\lvert 3-(-1)\rvert=4\)，恰等于两圆半径，所以该直线与两圆都相切. 故可填 \(x=-1\).
\end{solution}

\begin{problem}
若曲线 \(y=(x+a)\e^x\) 有两条过坐标原点的切线，则 \(a\) 的取值范围是\fillinblank{}.
\end{problem}

\begin{answer}
\((-\infty,-4)\cup(0,+\infty)\)
\end{answer}

\begin{solution}
设切点横坐标为 \(x_0\)，记 \(f(x)=(x+a)\e^x\)，则 \(f'(x)=(x+a+1)\e^x\). 切线方程为 \(y-(x_0+a)\e^{x_0}=(x_0+a+1)\e^{x_0}(x-x_0)\). 因为切线过原点，所以 \(-(x_0+a)\e^{x_0}=-(x_0+a+1)\e^{x_0}x_0\). 又 \(\e^{x_0}\ne 0\)，约去后得 \(x_0+a=(x_0+a+1)x_0\)，整理为 \(x_0^2+ax_0-a=0\). 若 \(x_0=1\)，上式化为 \(1=0\)，不可能，所以切点横坐标满足 \(x_0\ne 1\). 由 \(x_0+a=(x_0+a+1)x_0\) 得 \(a=\frac{x_0^2}{1-x_0}\)，进而 \(x_0+a+1=\frac{1}{1-x_0}\)，故切线斜率可记为 \(m(x_0)=\frac{\e^{x_0}}{1-x_0}\). 该二次方程的判别式为 \(\Delta=a^2+4a\). 当 \(\Delta>0\)，即 \(a<-4\) 或 \(a>0\) 时，方程有两个不同实根 \(r_1\ne r_2\)，且乘积为 \(-a\ne 0\)，故两根均非零，每根给出一条过原点的切线. 下证两根给出的切线不同. 当 \(a>0\) 时，两根一正一负；正根 \(r=\frac{-a+\sqrt{a^2+4a}}{2}<1\)，因为 \(\sqrt{a^2+4a}<a+2\)，故两根皆小于 \(1\). 而 \(m'(r)=\frac{(2-r)\e^r}{(1-r)^2}>0\)（\(r<1\)），所以 \(m\) 在 \((-\infty,1)\) 上严格递增，两根斜率不同. 当 \(a<-4\) 时，两根之和与乘积均为 \(-a>4\)，故两根皆大于 \(1\)；由 \((r_1-1)(r_2-1)=r_1r_2-(r_1+r_2)+1=1\)，记 \(t_1=r_1-1\)，\(t_2=r_2-1\)，则 \(t_1t_2=1\) 且 \(t_1\ne 1\). 若两切线斜率相等，则 \(\frac{\e^{r_1}}{1-r_1}=\frac{\e^{r_2}}{1-r_2}\)，即 \(\frac{\e^{t_1}}{t_1}=\frac{\e^{t_2}}{t_2}\)，结合 \(t_2=\frac{1}{t_1}\) 得 \(t_1-\frac{1}{t_1}=2\ln t_1\). 记 \(\eta(t)=t-\frac1t-2\ln t\)（\(t>0\)，\(t\ne 1\)），则 \(\eta'(t)=\left(1-\frac1t\right)^2>0\)，所以 \(\eta\) 在 \((0,1)\) 和 \((1,+\infty)\) 上分别严格递增，又 \(\eta(1)=0\)，故 \(t-\frac1t=2\ln t\) 只有解 \(t=1\)，与 \(t_1\ne 1\) 矛盾. 因此两根斜率不同，给出两条不同切线. 反之当 \(\Delta\le 0\) 时二次方程至多一根，至多一条过原点的切线. 故取值范围是 \((-\infty,-4)\cup(0,+\infty)\).
\end{solution}

\begin{problem}
已知椭圆 \(C:\ \frac{x^2}{a^2}+\frac{y^2}{b^2}=1\)（\(a>b>0\)），\(C\) 的上顶点为 \(A\)，两个焦点为 \(F_1\)，\(F_2\)，离心率为 \(\frac12\). 过 \(F_1\) 且垂直于 \(AF_2\) 的直线与 \(C\) 交于 \(D\)，\(E\) 两点，\(|DE|=6\)，则 \(\triangle ADE\) 的周长是\fillinblank{}.
\end{problem}

\begin{answer}
\(13\)
\end{answer}

\begin{solution}
因为离心率 \(e=\frac{c}{a}=\frac{1}{2}\)，所以 \(a=2c\)，\(b^2=a^2-c^2=3c^2\)，椭圆方程为 \(\frac{x^2}{4c^2}+\frac{y^2}{3c^2}=1\). 记上顶点为 \(A(0,b)\)，焦点为 \(F_1(-c,0)\)，\(F_2(c,0)\). 因为 \(\lvert AF_1\rvert=\lvert AF_2\rvert=a\) 且 \(\lvert F_1F_2\rvert=2c=a\)，所以 \(\triangle AF_1F_2\) 为边长 \(a\) 的等边三角形. 过 \(F_1\) 且垂直于 \(AF_2\) 的直线即等边三角形从 \(F_1\) 向对边 \(AF_2\) 作的高，也是 \(AF_2\) 的垂直平分线，故其上任一点到 \(A\) 与到 \(F_2\) 距离相等，特别有 \(\lvert AD\rvert=\lvert DF_2\rvert\)，\(\lvert AE\rvert=\lvert EF_2\rvert\). 因为 \(AF_2\) 斜率为 \(-\frac{b}{c}=-\sqrt{3}\)，所以 \(DE\) 斜率为 \(\frac{1}{\sqrt{3}}\)，方程为 \(y=\frac{1}{\sqrt{3}}(x+c)\). 联立椭圆方程得 \(\frac{x^2}{4c^2}+\frac{(x+c)^2}{9c^2}=1\)，即 \(13x^2+8cx-32c^2=0\). 设 \(D\)，\(E\) 横坐标为 \(x_D\)，\(x_E\)，则 \(x_D+x_E=-\frac{8c}{13}\)，\(x_Dx_E=-\frac{32c^2}{13}\). 由弦长公式得 \(\lvert DE\rvert=\sqrt{1+\frac{1}{3}}\lvert x_D-x_E\rvert=\frac{2}{\sqrt{3}}\sqrt{(x_D+x_E)^2-4x_Dx_E}=\frac{2}{\sqrt{3}}\sqrt{\frac{64c^2}{169}+\frac{128c^2}{13}}=\frac{48}{13}c\). 由 \(\lvert DE\rvert=6\) 得 \(c=\frac{13}{8}\)，从而 \(a=2c=\frac{13}{4}\). 因为 \(F_1\) 在椭圆内，\(D\)，\(E\) 分居 \(F_1\) 两侧，所以 \(\lvert DE\rvert=\lvert DF_1\rvert+\lvert EF_1\rvert\). 于是 \(\triangle ADE\) 周长等于 \((\lvert AD\rvert+\lvert AE\rvert)+\lvert DE\rvert=(\lvert DF_2\rvert+\lvert DF_1\rvert)+(\lvert EF_2\rvert+\lvert EF_1\rvert)=2a+2a=4a=13\). 故答案为 \(13\).
\end{solution}
\section{解答题}

\begin{problem}
记 \(S_n\) 为数列 \(\{a_n\}\) 的前 \(n\) 项和，已知 \(a_1=1\)，数列 \(\left\{\frac{S_n}{a_n}\right\}\) 是公差为 \(\frac13\) 的等差数列.

\begin{enumerate}
\item 求 \(\{a_n\}\) 的通项公式；

\item 证明：\(\frac1{a_1}+\frac1{a_2}+\cdots+\frac1{a_n}<2\).
\end{enumerate}
\end{problem}

\begin{solution}
\begin{enumerate}
\item 因为
\(\frac{S_1}{a_1}=1\)，且 \(\left\{\frac{S_n}{a_n}\right\}\) 是公差为 \(\frac13\) 的等差数列，所以
\(\frac{S_n}{a_n}=1+\frac{n-1}{3}=\frac{n+2}{3}\).
于是
\(S_n=\frac{n+2}{3}a_n\).
当 \(n\ge 2\) 时，\(S_{n-1}=\frac{n+1}{3}a_{n-1}\)，所以
\(a_n=S_n-S_{n-1}=\frac{n+2}{3}a_n-\frac{n+1}{3}a_{n-1}\)，整理得
\((n-1)a_n=(n+1)a_{n-1}\).
因为 \(n\ge 2\) 时 \(n-1>0\)，且由 \(a_1=1\) 递推，所以各 \(a_n\ne 0\)，所以
\(\frac{a_n}{a_{n-1}}=\frac{n+1}{n-1}\)（\(n\ge 2\)）.
因此
\(a_n=a_1\cdot \frac{a_2}{a_1}\cdot \frac{a_3}{a_2}\cdots \frac{a_n}{a_{n-1}}=1\cdot \frac31\cdot \frac42\cdots \frac{n}{n-2}\cdot \frac{n+1}{n-1}=\frac{n(n+1)}{2}\).
又当 \(n=1\) 时，\(\frac{1\cdot 2}{2}=1=a_1\)，所以对一切正整数 \(n\)，\(a_n=\frac{n(n+1)}{2}\).
\item 由（1）知
\(\frac1{a_n}=\frac{2}{n(n+1)}=2\left(\frac1n-\frac1{n+1}\right)\).
因此
\(\frac1{a_1}+\frac1{a_2}+\cdots+\frac1{a_n}=2\sum_{k=1}^n\left(\frac1k-\frac1{k+1}\right)=2\left(1-\frac1{n+1}\right)<2\).
证毕.
\end{enumerate}
\end{solution}

\begin{problem}
记 \(\triangle ABC\) 的内角 \(A\)，\(B\)，\(C\) 的对边分别为 \(a\)，\(b\)，\(c\)，已知 \(\frac{\cos A}{1+\sin A}=\frac{\sin 2B}{1+\cos 2B}\).

\begin{enumerate}
\item 若 \(C=\frac{2\pi}{3}\)，求 \(B\)；

\item 求 \(\frac{a^2+b^2}{c^2}\) 的最小值.
\end{enumerate}
\end{problem}

\begin{solution}
\begin{enumerate}
\item 因为 \(B\in(0,\pi)\) 且 \(1+\cos 2B=2\cos^2B\)，所以 \(B\ne \frac{\pi}{2}\)，且
\(\frac{\sin 2B}{1+\cos 2B}=\frac{2\sin B\cos B}{2\cos^2B}=\frac{\sin B}{\cos B}\).
所以已知条件化为
\(\frac{\cos A}{1+\sin A}=\frac{\sin B}{\cos B}\)，即
\(\cos A\cos B=\sin B(1+\sin A)\)，整理得
\(\cos A\cos B-\sin A\sin B=\sin B\)，即
\(\cos(A+B)=\sin B\).
因为 \(A+B=\pi-C\)，所以
\(\sin B=-\cos C\).
当 \(C=\frac{2\pi}{3}\) 时，
\(\sin B=-\cos\frac{2\pi}{3}=\frac12\).
又此时 \(A+B=\frac{\pi}{3}\)，所以 \(0<B<\frac{\pi}{3}\)，故
\(B=\frac{\pi}{6}\).
\item 由（1）已证 \(\sin B=-\cos C>0\)，因为 \(B\in(0,\pi)\)，\(C\in(0,\pi)\)，所以
\(\frac{\pi}{2}<C<\pi\)，\(0<B<\frac{\pi}{2}\).
又 \(-\cos C=\sin\left(C-\frac{\pi}{2}\right)\)，所以
\(\sin B=\sin\left(C-\frac{\pi}{2}\right)\).
因为 \(B\in\left(0,\frac{\pi}{2}\right)\)，\(C-\frac{\pi}{2}\in\left(0,\frac{\pi}{2}\right)\)，而正弦函数在 \(\left(0,\frac{\pi}{2}\right)\) 上单调递增，所以
\(B=C-\frac{\pi}{2}\)，即
\(C=\frac{\pi}{2}+B\).
从而 \(A=\pi-B-C=\frac{\pi}{2}-2B\).
因为 \(A>0\)，所以 \(B<\frac{\pi}{4}\)，结合 \(B>0\) 得
\(B\in\left(0,\frac{\pi}{4}\right)\)，\(C\in\left(\frac{\pi}{2},\frac{3\pi}{4}\right)\).
由正弦定理，
\(\frac{a^2+b^2}{c^2}=\frac{\sin^2A+\sin^2B}{\sin^2C}\).
因为 \(\sin A=\sin\left(\frac{\pi}{2}-2B\right)=\cos 2B\)，\(\sin C=\sin\left(\frac{\pi}{2}+B\right)=\cos B\)，所以
\(\frac{a^2+b^2}{c^2}=\frac{\cos^22B+\sin^2B}{\cos^2B}\).
又 \(\cos 2B=2\cos^2B-1\)，\(\sin^2B=1-\cos^2B\)，所以分子为 \(\left(2\cos^2B-1\right)^2+1-\cos^2B\)，即
\(\frac{a^2+b^2}{c^2}=4\cos^2B+\frac{2}{\cos^2B}-5\).
记 \(t=\cos^2B\)，因为 \(B\in\left(0,\frac{\pi}{4}\right)\)，所以 \(t\in\left(\frac12,1\right)\).
于是 \(\frac{a^2+b^2}{c^2}=4t+\frac{2}{t}-5\).
因为 \(4t>0\)，\(\frac{2}{t}>0\)，由均值不等式，
\(4t+\frac{2}{t}\ge 2\sqrt{4t\cdot \frac{2}{t}}=2\sqrt8=4\sqrt2\)，所以
\(\frac{a^2+b^2}{c^2}\ge 4\sqrt2-5\).
当且仅当 \(4t=\frac{2}{t}\)，即 \(t=\frac{\sqrt2}{2}\) 时取等号.
因为 \(\frac12<\frac{\sqrt2}{2}<1\)，故该 \(t\) 落在 \(\left(\frac12,1\right)\) 内可取，对应 \(B\in\left(0,\frac{\pi}{4}\right)\) 存在.
所以 \(\frac{a^2+b^2}{c^2}\) 的最小值为 \(4\sqrt2-5\).
\end{enumerate}
\end{solution}

\begin{problem}
如图，直三棱柱 \(ABC-A_1B_1C_1\) 的体积为 \(4\)，\(\triangle A_1BC\) 的面积为 \(2\sqrt2\).

\begin{enumerate}
\item 求 \(A\) 到平面 \(A_1BC\) 的距离；

\item 设 \(D\) 为 \(A_1C\) 的中点，\(AA_1=AB\)，平面 \(A_1BC\perp\) 平面 \(ABB_1A_1\)，求二面角 \(A-BD-C\) 的正弦值.
\end{enumerate}

\bitmapfigure[max width=0.28\textwidth]{img/2022/new_gaokao_paper_1/q19_fig1.png}
\end{problem}

\begin{solution}
\solutionfigure{\bitmapfigure[max width=0.27\textwidth]{img/2022/new_gaokao_paper_1/q19_solution_fig1.png}}
\begin{enumerate}
\item 设点 \(A\) 到平面 \(A_1BC\) 的距离为 \(h\)，则
\(V_{A-A_1BC}=\frac13S_{\triangle A_1BC}\cdot h=\frac13\cdot 2\sqrt2\cdot h=\frac{2\sqrt2}{3}h\).
又三棱锥 \(A-A_1BC\) 与三棱锥 \(A_1-ABC\) 为同一几何体，而
\(V_{A_1-ABC}=\frac13V_{ABC-A_1B_1C_1}=\frac43\)，所以
\(V_{A-A_1BC}=\frac43\).
解得 \(h=\sqrt2\).
所以点 \(A\) 到平面 \(A_1BC\) 的距离为 \(\sqrt2\).
\item 取 \(A_1B\) 的中点 \(E\)，连接 \(AE\)，因为 \(AA_1=AB\)，所以
\(AE\perp A_1B\).
又平面 \(A_1BC\perp\) 平面 \(ABB_1A_1\)，交线为 \(A_1B\)，且 \(AE\subset\) 平面 \(ABB_1A_1\)，所以
\(AE\perp\) 平面 \(A_1BC\).
因此 \(AE\) 即为点 \(A\) 到平面 \(A_1BC\) 的距离，由（1）得 \(AE=\sqrt2\).
在直三棱柱中，\(BB_1\perp\) 平面 \(ABC\)，因为 \(BC\subset\) 平面 \(ABC\)，所以 \(BB_1\perp BC\).
又 \(AE\perp\) 平面 \(A_1BC\)，\(BC\subset\) 平面 \(A_1BC\)，所以 \(AE\perp BC\).
因为 \(AE\)，\(BB_1\subset\) 平面 \(ABB_1A_1\) 且相交，所以
\(BC\perp\) 平面 \(ABB_1A_1\).
因此 \(BC\perp AB\)，\(BC\perp BB_1\)，而 \(AB\perp BB_1\)，故 \(BC\)，\(BA\)，\(BB_1\) 两两垂直.
以 \(B\) 为原点，分别以 \(BC\)，\(BA\)，\(BB_1\) 为 \(x\) 轴、\(y\) 轴、\(z\) 轴建立空间直角坐标系.
因为 \(ABB_1A_1\) 为 \(AA_1=AB\) 的矩形，故为正方形，所以 \(A_1B=\sqrt2AB\)，\(AE=\frac12A_1B=\frac{\sqrt2}{2}AB\).
由 \(AE=\sqrt2\) 得 \(AB=2\)，所以 \(AA_1=2\)，\(A_1B=2\sqrt2\).
又 \(BC\perp\) 平面 \(ABB_1A_1\)，所以 \(BC\perp A_1B\)，故 \(S_{\triangle A_1BC}=\frac12A_1B\cdot BC=2\sqrt2\)，解得 \(BC=2\).
则 \(A(0,2,0)\)，\(A_1(0,2,2)\)，\(B(0,0,0)\)，\(C(2,0,0)\)，所以 \(A_1C\) 的中点 \(D(1,1,1)\).
则 \(\overrightarrow{BD}=(1,1,1)\)，\(\overrightarrow{BA}=(0,2,0)\)，\(\overrightarrow{BC}=(2,0,0)\).
设平面 \(ABD\) 的一个法向量 \(\bs m=(x,y,z)\)，则由 \(\bs m\cdot \overrightarrow{BD}=x+y+z=0\) 和 \(\bs m\cdot \overrightarrow{BA}=2y=0\) 可取 \(\bs m=(1,0,-1)\).
设平面 \(BDC\) 的一个法向量 \(\bs n=(a,b,c)\)，则由 \(\bs n\cdot \overrightarrow{BD}=a+b+c=0\) 和 \(\bs n\cdot \overrightarrow{BC}=2a=0\) 可取 \(\bs n=(0,1,-1)\).
则 \(\cos\langle\bs m,\bs n\rangle=\frac{\bs m\cdot \bs n}{|\bs m||\bs n|}=\frac{1}{\sqrt2\times \sqrt2}=\frac12\)，所以二面角 \(A-BD-C\) 的正弦值为 \(\sqrt{1-\left(\frac12\right)^2}=\frac{\sqrt3}{2}\).
\end{enumerate}
\end{solution}

\begin{problem}
一医疗团队为研究某地的一种地方性疾病与当地居民的卫生习惯（卫生习惯分为良好和不够良好两类）的关系，在已患该疾病的病例中随机调查了 \(100\) 例（称为病例组），同时在未患该疾病的人群中随机调查了 \(100\) 人（称为对照组），得到如下数据：

\begin{center}
\begin{tabular}{|c|c|c|}
\hline
 & 不够良好 & 良好 \\
\hline
病例组 & 40 & 60 \\
\hline
对照组 & 10 & 90 \\
\hline
\end{tabular}
\end{center}

\begin{enumerate}
\item 能否有 \(99\%\) 的把握认为患该疾病群体与未患该疾病群体的卫生习惯有差异？

\item 从该地的人群中任选一人，\(A\) 表示事件“选到的人卫生习惯不够良好”，\(B\) 表示事件“选到的人患有该疾病”. \(\frac{P(B\mid A)}{P(\overline B\mid A)}\) 与 \(\frac{P(B\mid \overline A)}{P(\overline B\mid \overline A)}\) 的比值是卫生习惯不够良好对患该疾病风险程度的一项度量指标，记该指标为 \(R\).

\begin{enumerate}
\item 证明：\(R=\frac{P(B\mid A)}{P(\overline B\mid A)}\div \frac{P(B\mid \overline A)}{P(\overline B\mid \overline A)}=\frac{P(A\mid B)}{P(\overline A\mid B)}\cdot \frac{P(\overline A\mid \overline B)}{P(A\mid \overline B)};\)

\item 利用该调查数据，给出 \(P(A\mid B)\)，\(P(A\mid \overline B)\) 的估计值，并利用（1）的结果给出 \(R\) 的估计值.
\end{enumerate}

附：\(K^2=\frac{n(ad-bc)^2}{(a+b)(c+d)(a+c)(b+d)}\)，

\begin{center}
\begin{tabular}{|c|c|c|c|}
\hline
\(P(K^2\ge k)\) & 0.050 & 0.010 & 0.001 \\
\hline
\(k\) & 3.841 & 6.635 & 10.828 \\
\hline
\end{tabular}
\end{center}
\end{enumerate}
\end{problem}

\begin{solution}
\begin{enumerate}
\item 补全列联表：
{\centering
\begin{tabular}{|c|c|c|c|}
\hline
& 不够良好 & 良好 & 合计 \\
\hline
病例组 & 40 & 60 & 100 \\
\hline
对照组 & 10 & 90 & 100 \\
\hline
合计 & 50 & 150 & 200 \\
\hline
\end{tabular}
\par}
因为 \(40\times 90-10\times 60=3000\)，所以
\(K^2=\frac{200(40\times 90-10\times 60)^2}{100\times 100\times 50\times 150}=24\).
因为 \(24>6.635\)，所以有 \(99\%\) 的把握认为两类人群的卫生习惯有差异.
\item \begin{enumerate}
\item 证明：
\(R=\frac{P(B\mid A)}{P(\overline B\mid A)}\div \frac{P(B\mid \overline A)}{P(\overline B\mid \overline A)}=\frac{\frac{P(AB)}{P(A)}}{\frac{P(A\overline B)}{P(A)}}\div \frac{\frac{P(\overline AB)}{P(\overline A)}}{\frac{P(\overline A\overline B)}{P(\overline A)}}=\frac{P(AB)}{P(A\overline B)}\cdot \frac{P(\overline A\overline B)}{P(\overline AB)}=\frac{P(A\mid B)}{P(\overline A\mid B)}\cdot \frac{P(\overline A\mid \overline B)}{P(A\mid \overline B)}\).
\item 由调查数据估计 \(P(A\mid B)=\frac{40}{100}=\frac25\)，\(P(A\mid \overline B)=\frac{10}{100}=\frac1{10}\).
于是 \(P(\overline A\mid B)=\frac35\)，\(P(\overline A\mid \overline B)=\frac9{10}\).
代入（1）得 \(R=\frac{\frac25}{\frac35}\cdot \frac{\frac9{10}}{\frac1{10}}=6\).
\end{enumerate}
\end{enumerate}
\end{solution}

\begin{problem}
已知点 \(A(2,1)\) 在双曲线 \(C:\ \frac{x^2}{a^2}-\frac{y^2}{a^2-1}=1\)（\(a>1\)） 上，直线 \(l\) 交 \(C\) 于 \(P\)，\(Q\) 两点，直线 \(AP\)，\(AQ\) 的斜率之和为 \(0\).

\begin{enumerate}
\item 求 \(l\) 的斜率；

\item 若 \(\tan \angle PAQ=2\sqrt2\)，求 \(\triangle PAQ\) 的面积.
\end{enumerate}
\end{problem}

\begin{solution}
\begin{enumerate}
\item 因为 \(A(2,1)\in C\)，所以 \(\frac{4}{a^2}-\frac{1}{a^2-1}=1\)，整理得 \(a^4-4a^2+4=0\)，即 \(\left(a^2-2\right)^2=0\)，解得 \(a^2=2\).
故双曲线方程为 \(\frac{x^2}{2}-y^2=1\).
若直线 \(l\) 的斜率不存在，则可设 \(l:x=t\).
当 \(t=2\) 时，直线 \(l\) 经过点 \(A\)，此时直线 \(AP\) 或 \(AQ\) 的斜率不存在，与题意不符.
当 \(t\ne 2\) 时，直线 \(l\) 与双曲线 \(C\) 的交点可设为 \(P(t,y_1)\)，\(Q(t,y_2)\)，由 \(\frac{t^2}{2}-y^2=1\) ，所以 \(y_1+y_2=0\)，从而 \(k_{AP}+k_{AQ}=\frac{y_1-1}{t-2}+\frac{y_2-1}{t-2}=\frac{y_1+y_2-2}{t-2}=-\frac{2}{t-2}\ne 0\)，这与已知条件矛盾.
所以直线 \(l\) 的斜率存在，设 \(l:y=kx+m\)，\(P(x_1,y_1)\)，\(Q(x_2,y_2)\).
联立 \(y=kx+m\) 和 \(\frac{x^2}{2}-y^2=1\)，得到 \(\left(1-2k^2\right)x^2-4mkx-2m^2-2=0\).
因为是一元二次方程，所以 \(k\ne \pm \frac{\sqrt2}{2}\)，且判别式 \(\Delta=16m^2k^2+4\left(1-2k^2\right)\left(2m^2+2\right)=8\left(m^2+1-2k^2\right)>0\)，即 \(m^2+1-2k^2>0\).
由韦达定理，\(x_1+x_2=\frac{4mk}{1-2k^2}=-\frac{4mk}{2k^2-1}\)，\(x_1x_2=\frac{-2m^2-2}{1-2k^2}=\frac{2m^2+2}{2k^2-1}\).
因为直线 \(AP\)，\(AQ\) 斜率存在，所以 \(P\)，\(Q\) 均不为 \(A\)，故 \(l\) 不过点 \(A\).
由 \(k_{AP}+k_{AQ}=0\) 得 \(\frac{y_1-1}{x_1-2}+\frac{y_2-1}{x_2-2}=0\)，即 \(\left(x_1-2\right)\left(kx_2+m-1\right)+\left(x_2-2\right)\left(kx_1+m-1\right)=0\)，即 \(2kx_1x_2+(m-1-2k)\left(x_1+x_2\right)-4(m-1)=0\).
代入韦达结果，通分后分子为 \(2k\left(2m^2+2\right)+(m-1-2k)(-4mk)-4(m-1)\left(2k^2-1\right)\).
展开后 \(4km^2\) 与 \(-4m^2k\) 相消，\(8mk^2\) 与 \(-8mk^2\) 相消，余下 \(4(k+mk+m+2k^2-1)=4(k+1)(m+2k-1)\)，所以化简得 \((k+1)(m+2k-1)=0\).
若 \(m=1-2k\)，则直线 \(l\) 过点 \(A\)，与题意不符，故 \(k=-1\).
\item 由（1）知 \(k=-1\)，则 \(x_1x_2=2m^2+2>0\)，所以 \(x_1\)，\(x_2\) 同号.
不妨设直线 \(PA\)，\(AQ\) 的倾斜角为 \(\alpha\)，\(\beta\)（\(\alpha<\frac{\pi}{2}<\beta\)），因为 \(k_{AP}+k_{AQ}=0\)，所以 \(\alpha+\beta=\pi\).
当 \(P\)，\(Q\) 均在双曲线左支时，\(\angle PAQ=2\alpha\)，所以 \(\tan 2\alpha=2\sqrt2\)，即 \(\sqrt2\tan^2\alpha+\tan\alpha-\sqrt2=0\)，分解得 \(\left(\sqrt2\tan\alpha-1\right)(\tan\alpha+\sqrt2)=0\)，解得 \(\tan\alpha=\frac{\sqrt2}{2}\)（负值舍去）.
此时 \(PA\) 的斜率等于渐近线 \(y=\frac{\sqrt2}{2}x\) 的斜率，与双曲线左支无交点，舍去.
当 \(P\)，\(Q\) 均在双曲线右支时，\(\tan\angle PAQ=2\sqrt2\)，所以 \(\tan(\beta-\alpha)=2\sqrt2\)，即 \(\tan 2\alpha=-2\sqrt2\)，也就是 \(\sqrt2\tan^2\alpha-\tan\alpha-\sqrt2=0\)，分解得 \(\left(\sqrt2\tan\alpha+1\right)(\tan\alpha-\sqrt2)=0\)，解得 \(\tan\alpha=\sqrt2\)（负值舍去）.
于是 \(PA:y=\sqrt2(x-2)+1\)，\(AQ:y=-\sqrt2(x-2)+1\).
联立 \(y=\sqrt2(x-2)+1\) 和 \(\frac{x^2}{2}-y^2=1\)，得到 \(\frac32x^2+2\left(\sqrt2-4\right)x+10-4\sqrt2=0\).
因为方程有一个根为 \(2\)，由韦达定理两根之和为 \(\frac{16-4\sqrt2}{3}\)，所以另一根为 \(\frac{10-4\sqrt2}{3}\)，得 \(P\left(\frac{10-4\sqrt2}{3},\frac{4\sqrt2-5}{3}\right)\). 联立 \(y=-\sqrt2(x-2)+1\) 和 \(\dfrac{x^2}{2}-y^2=1\)，得 \(\frac32x^2-2\left(\sqrt2+4\right)x+10+4\sqrt2=0\). 该方程同样以 \(2\) 为一根，由韦达定理两根之和为 \(\frac{16+4\sqrt2}{3}\)，所以另一根为 \(\frac{10+4\sqrt2}{3}\)，得 \(Q\left(\frac{10+4\sqrt2}{3},\frac{-4\sqrt2-5}{3}\right)\).
所以 \(PQ:x+y-\frac53=0\)，\(|PQ|=\frac{16}{3}\).
点 \(A\) 到直线 \(PQ\) 的距离为 \(d=\frac{\left|2+1-\frac53\right|}{\sqrt2}=\frac{2\sqrt2}{3}\).
故 \(S_{\triangle PAQ}=\frac12\cdot \frac{16}{3}\cdot \frac{2\sqrt2}{3}=\frac{16\sqrt2}{9}\).
\end{enumerate}
\end{solution}

\begin{problem}
已知函数 \(f(x)=\e^x-ax\)，\(g(x)=ax-\ln x\) 有相同的最小值.

\begin{enumerate}
\item 求 \(a\)；

\item 证明：存在直线 \(y=b\)，其与两条曲线 \(y=f(x)\) 和 \(y=g(x)\) 共有三个不同的交点，并且从左到右的三个交点的横坐标成等差数列.
\end{enumerate}
\end{problem}

\begin{solution}
\begin{enumerate}
\item 对 \(f(x)=\e^x-ax\)（\(x\in\R\)）有 \(f^{\prime}(x)=\e^x-a\).
当 \(a\le 0\) 时，\(f^{\prime}(x)>0\) 恒成立，所以 \(f\) 在 \(\R\) 上单调递增，无最小值，故 \(a>0\).
此时由 \(f^{\prime}(x)=0\) 得 \(x=\ln a\)，当 \(x<\ln a\) 时 \(f^{\prime}(x)<0\)，当 \(x>\ln a\) 时 \(f^{\prime}(x)>0\)，所以 \(f\) 的最小值为 \(f(\ln a)=a-a\ln a\).
对 \(g(x)=ax-\ln x\)（\(x>0\)）有 \(g^{\prime}(x)=a-\frac1x\).
因为 \(a>0\)，当 \(0<x<\frac1a\) 时 \(g^{\prime}(x)<0\)，故 \(g\) 在 \(\left(0,\frac1a\right)\) 上为减函数.
当 \(x>\frac1a\) 时 \(g^{\prime}(x)>0\)，故 \(g\) 在 \(\left(\frac1a,+\infty\right)\) 上为增函数.
故 \(g\) 的最小值为 \(g\left(\frac1a\right)=1-\ln\frac1a=1+\ln a\).
因为 \(f\) 和 \(g\) 有相同的最小值，故 \(a-a\ln a=1+\ln a\)，整理为 \(a-1=(a+1)\ln a\)，即 \(\frac{a-1}{a+1}=\ln a\).
设 \(\phi(a)=\frac{a-1}{a+1}-\ln a\)（\(a>0\)），则 \(\phi^{\prime}(a)=\frac{2}{(1+a)^2}-\frac1a=-\frac{a^2+1}{a(1+a)^2}<0\).
故 \(\phi\) 单调递减，而 \(\phi(1)=0\)，所以唯一解为 \(a=1\).
\item 由（1）知 \(f(x)=\e^x-x\)，\(g(x)=x-\ln x\)，最小值均为 \(1\).
当 \(b>1\) 时，设 \(S(x)=\e^x-x-b\)，则 \(S^{\prime}(x)=\e^x-1\).
当 \(x<0\) 时 \(S^{\prime}(x)<0\)，当 \(x>0\) 时 \(S^{\prime}(x)>0\)，故 \(S\) 在 \((-\infty,0)\) 上为减函数，在 \((0,+\infty)\) 上为增函数，所以 \(S(x)_{\min}=S(0)=1-b<0\).
而 \(S(-b)=\e^{-b}>0\)，\(S(b)=\e^b-2b\).
设 \(u(b)=\e^b-2b\)（\(b>1\)），则 \(u^{\prime}(b)=\e^b-2>\e-2>0\)，故 \(u\) 在 \((1,+\infty)\) 上为增函数，所以 \(u(b)>u(1)=\e-2>0\).
故 \(S(b)>0\)，所以 \(S(x)=0\) 有两个不同零点，即 \(\e^x-x=b\) 有两个不同解.
设 \(T(x)=x-\ln x-b\)（\(x>0\)），则 \(T^{\prime}(x)=\frac{x-1}{x}\).
当 \(0<x<1\) 时 \(T^{\prime}(x)<0\)，当 \(x>1\) 时 \(T^{\prime}(x)>0\)，故 \(T\) 在 \((0,1)\) 上为减函数，在 \((1,+\infty)\) 上为增函数，所以 \(T(x)_{\min}=T(1)=1-b<0\).
而 \(T(\e^{-b})=\e^{-b}>0\)，\(T(\e^b)=\e^b-2b>0\)，故 \(T(x)=0\) 有两个不同零点，即 \(x-\ln x=b\) 有两个不同解.
当 \(b=1\) 时，由（1）的单调性，所以两方程各仅有一个解，当 \(b<1\) 时均无解.
故若直线 \(y=b\) 与两曲线共有三个不同交点，则 \(b>1\).
设 \(h(x)=\e^x+\ln x-2x\)（\(x>0\)），则 \(h^{\prime}(x)=\e^x+\frac1x-2\).
设 \(s(x)=\e^x-x-1\)（\(x>0\)），则 \(s^{\prime}(x)=\e^x-1>0\)，故 \(s\) 在 \((0,+\infty)\) 上为增函数，所以 \(s(x)>s(0)=0\)，即 \(\e^x>x+1\).
所以 \(h^{\prime}(x)>x+\frac1x-1\ge 2-1>0\)，故 \(h\) 在 \((0,+\infty)\) 上为增函数.
而 \(h(1)=\e-2>0\)，并且 \(h\left(\frac1{\e^3}\right)=\e^{\frac1{\e^3}}-3-\frac2{\e^3}<\e-3-\frac2{\e^3}<0\)，故 \(h\) 在 \((0,+\infty)\) 上有且只有一个零点 \(x_0\)，且 \(\frac1{\e^3}<x_0<1\).
当 \(0<x<x_0\) 时 \(h(x)<0\)，即 \(\e^x-x<x-\ln x\)，即 \(f(x)<g(x)\).
当 \(x>x_0\) 时 \(h(x)>0\)，即 \(f(x)>g(x)\).
因此若直线 \(y=b\) 与两曲线共有三个不同交点，则 \(b=f(x_0)=g(x_0)>1\).
此时 \(\e^x-x=b\) 有两个不同根 \(x_1\)，\(x_0\)（\(x_1<0<x_0\)），\(x-\ln x=b\) 有两个不同根 \(x_0\)，\(x_4\)（\(0<x_0<1<x_4\)）.
故 \(\e^{x_1}-x_1=b\)，\(\e^{x_0}-x_0=b\)，\(x_4-\ln x_4-b=0\)，\(x_0-\ln x_0-b=0\).
所以 \(x_4-b=\ln x_4\)，即 \(\e^{x_4-b}=x_4\)，从而 \(\e^{x_4-b}-\left(x_4-b\right)-b=0\)，故 \(x_4-b\) 为方程 \(\e^x-x=b\) 的解.
又由 \(x_0-\ln x_0-b=0\) 得 \(\ln x_0=x_0-b\)，即 \(\e^{x_0-b}=x_0\)，从而 \(\e^{x_0-b}-\left(x_0-b\right)-b=0\)，故 \(x_0-b\) 也为方程 \(\e^x-x=b\) 的解.
又 \(\e^{x_1}-x_1=b\) 可化为 \(\e^{x_1}=x_1+b\)，即 \(x_1-\ln\left(x_1+b\right)=0\)，即 \(\left(x_1+b\right)-\ln\left(x_1+b\right)-b=0\)，故 \(x_1+b\) 为方程 \(x-\ln x=b\) 的解.
又 \(\e^{x_0}-x_0=b\) 可化为 \(\e^{x_0}=x_0+b\)，即 \(x_0-\ln\left(x_0+b\right)=0\)，即 \(\left(x_0+b\right)-\ln\left(x_0+b\right)-b=0\)，故 \(x_0+b\) 也为方程 \(x-\ln x=b\) 的解.
所以 \(\{x_1,x_0\}=\{x_0-b,x_4-b\}\)，而 \(b>1\)，\(x_0\in(0,1)\)，故 \(x_0-b<0<x_0\)，因此 \(x_0-b\) 必为负根 \(x_1\)，\(x_4-b\) 必为正根 \(x_0\)，即 \(x_1=x_0-b\)，\(x_0=x_4-b\)，于是 \(x_1+x_4=2x_0\).
所以直线 \(y=b\) 与两条曲线从左到右的三个交点的横坐标成等差数列.
证毕.
\end{enumerate}
\end{solution}
